% \documentclass[IB]{TFUOC}%IB: CASTELLANO, CAT: CATALÁN, ENG: ANGLÈS

% --- AÑADE ESTO AQUÍ ---
%\usepackage{amsmath, amssymb}
%\usepackage{hyperref}
%\newcommand{\citeonline}[1]{\cite{#1}}
%\setlength{\headheight}{43pt}
%\usepackage{comment}
%\setlength{\parskip}{0.6em}


%\usepackage[numbers]{natbib}

%\usepackage{geometry} % Si no está ya cargado por la clase
%\setlength{\headheight}{43pt} % El valor que te pedía el log (redondeado)
%\addtolength{\topmargin}{-8pt} % Compensa el movimiento para que no baje el texto

% En el preámbulo debes añadir: 
%\usepackage{pgfgantt} 

% Ajuste de cabecera solicitado por MiKTeX
%\setlength{\headheight}{43pt}

%\documentclass[IB]{TFUOC}%IB: CASTELLANO, CAT: CATALÁN, ENG: ANGLÈS

%\usepackage{amsmath, amssymb}
%\usepackage[numbers,sort&compress]{natbib}
%\usepackage[colorlinks=true,linkcolor=black,citecolor=black,urlcolor=blue]{hyperref}
%\usepackage{comment}
%\usepackage{geometry}
%\usepackage{pgfgantt}

%\setlength{\headheight}{43pt}
%\addtolength{\topmargin}{-8pt}
%\setlength{\parskip}{0.6em}

\documentclass[IB]{TFUOC}

\usepackage{amsmath, amssymb}
\usepackage[numbers,sort&compress]{natbib}
\usepackage[colorlinks=true,linkcolor=black,citecolor=black,urlcolor=blue]{hyperref}
\usepackage{comment}
\usepackage{geometry}
\usepackage{pgfgantt}

\setlength{\headheight}{43pt}
\addtolength{\topmargin}{-8pt}
\setlength{\parskip}{0.6em}

%Introducción de datos del trabajo
\title{Optimización Dinámica de Carteras Híbridas de ETFs mediante Deep Reinforcement Learning: Integrando Criptoactivos en la Inversión Tradicional.}
\titcrt{Deep Reinforcement Learning} %Título corto que aparecerá a la cabecera
\author{Marcos Marqués Primo}
\date{Febrero 2026}


\nomPDC{Ruben Perez Ibañez}
\nomPRA{Nombre Profesorado Responsable}
\titulac{Máster en Cienca de Datos}
\area{Área 4}
\idioma{Castellano}
\credits{15}
\parcla{palabra, clave, trabajo}

\licenc{ccBy}
%Posibles licencias
%ccByNcNd
%ccByNcSa
%ccByNc
%ccByNd
%ccBySa
%ccBy
%GNU
%copyright


%%%%%%%%%%
% Resumen en el idioma

\abstractidioma{
Este Trabajo Final de Máster aborda el diseño e implementación de un sistema de gestión dinámica de carteras híbridas basado en \textit{Deep Reinforcement Learning} (DRL). El sistema combina activos tradicionales (renta variable y renta fija) con ETFs regulados sobre criptoactivos (IBIT y ETHA), permitiendo evaluar la integración de la nueva clase de activos digitales en la inversión institucional. La política de asignación se modela mediante un agente \textit{Proximal Policy Optimization} (PPO) entrenado sobre un entorno Gymnasium personalizado, en el que el espacio de estados integra indicadores técnicos y métricas de riesgo, y la recompensa penaliza explícitamente el \textit{Maximum Drawdown} y el \textit{turnover}. La validación se realiza mediante un protocolo \textit{out-of-sample} estricto con esquemas \textit{walk-forward} y \textit{expanding window}, y se compara el desempeño del agente frente a \textit{baselines} financieros clásicos (60/40, \textit{Equal-Weight}, \textit{Buy \& Hold} y Markowitz). Los resultados permiten cuantificar en qué medida un agente de aprendizaje por refuerzo profundo aporta una mejora reproducible en términos de rentabilidad ajustada al riesgo respecto a estrategias estáticas de referencia.
}

% Resumen en inglés.
\abstractenglish{

}

\begin{document}
\estructura

\tableofcontents

\listoffigures

\listoftables



\chapter{Descripción y justificación}
\section{Descripción y justificación}

\subsection{Interés y Relevancia}
La integración de activos cripto en carteras institucionales a través de ETFs ha marcado un hito en la eficiencia de los mercados. Sin embargo, la volatilidad inherente de estos activos requiere una gestión más activa y sofisticada que la tradicional. Este proyecto justifica su relevancia al proponer una solución tecnológica que utiliza la IA para mitigar riesgos y capitalizar las oportunidades de crecimiento de los activos digitales dentro de un marco de inversión diversificado y regulado.

\section{Contexto y justificación del trabajo}

La aprobación en enero de 2024 de los primeros ETFs \textit{spot} sobre Bitcoin en Estados Unidos, y muy particularmente del \textit{iShares Bitcoin Trust} (IBIT), supone un cambio cualitativo en el acceso institucional a los criptoactivos. Por primera vez, los gestores patrimoniales pueden construir exposición a Bitcoin a través de un instrumento regulado, custodiado por un \textit{prime broker} convencional y negociado en mercados organizados, eliminando una parte sustancial de las fricciones operativas y reputacionales que limitaban su uso en carteras institucionales. Esta evolución abre un escenario en el que la integración de activos digitales en carteras tradicionales deja de ser una cuestión hipotética y se convierte en un problema práctico de asignación de activos.

Sin embargo, los marcos clásicos de optimización de carteras —en particular la teoría moderna de Markowitz— se basan en supuestos (normalidad de los retornos, estacionariedad de las correlaciones, ausencia de fricciones) que no se sostienen cuando se incorporan activos con la dinámica observada en los criptoactivos: colas pesadas, \textit{volatility clustering} y dependencia condicional con el mercado tradicional que se intensifica precisamente en los episodios de tensión, justo cuando la diversificación es más necesaria. Una gestión basada en parámetros estimados sobre ventanas históricas y reasignaciones puntuales resulta, en este contexto, especialmente vulnerable.

El presente trabajo se justifica, por tanto, en la confluencia de tres factores: (i) la disponibilidad de un instrumento regulado que permite operar criptoactivos como un ETF cualquiera, (ii) la inadecuación de los marcos estáticos clásicos para capturar la dinámica de estos activos, y (iii) la madurez alcanzada por los algoritmos de aprendizaje por refuerzo profundo —y en particular por \textit{Proximal Policy Optimization}— como herramientas para abordar problemas de decisión secuencial bajo incertidumbre con espacios de acción continuos. La intersección de estos tres factores delimita un problema relevante, abordable y poco explorado: cómo aprender una política dinámica de asignación que permita integrar criptoactivos en carteras híbridas controlando explícitamente el riesgo y los costes de transacción.

\section{Definición de los objetivos}
\subsection{Objetivo Principal}
\begin{itemize}
    \item Desarrollar un agente de DRL que optimice la asignación de capital en una cartera híbrida de ETFs, buscando mejorar el Ratio de Sharpe/Sortino frente a estrategias de referencia, incorporando el control de drawdown y costes. 
\end{itemize}


\subsection{Objetivos Secundarios}
\begin{itemize}
    \item \textbf{Benchmarking Comparativo:} Evaluar el rendimiento del agente frente a cuatro \textit{baselines} financieros: 1) Cartera \textit{Equal-Weight} con rebalanceo mensual, 2) Cartera clásica 60/40 (Acciones/Bonos), 3) Estrategia \textit{Buy \& Hold}, y 4) Modelo de Media-Varianza de Markowitz con una ventana de estimación  a 12 meses.
    \item \textbf{Optimización de la Función de Recompensa:} Diseñar una función de recompensa compuesta que penalice explícitamente el \textit{Maximum Drawdown} (MDD), los costes de transacción (comisiones y \textit{slippage}) y el \textit{Turnover} excesivo, buscando posibilitar una estrategia financiera estable.
    \item \textbf{Validación de Robustez ante Cambios de Régimen:} Analizar la capacidad de adaptación del agente mediante un protocolo de evaluación \textit{Out-of-Sample} estricto, realizando pruebas específicas en periodos de alta volatilidad y cambios de régimen (ej. correcciones bruscas en el mercado cripto vs. estabilidad en renta fija).
\end{itemize}

% --- SECCIÓN DE METODOLOGÍA ---

\section{Enfoque Metodológico General}
\label{sec:enfoque_metodologico}

El trabajo adopta como marco general la \textit{Design Science Research} (DSR), orientada a la construcción y evaluación de un artefacto tecnológico que resuelve un problema práctico relevante. El artefacto es un sistema integral de gestión dinámica de carteras híbridas que articula cinco fases sucesivas:

\begin{enumerate}
    \item \textbf{Adquisición de datos e ingeniería de características}, sobre un universo híbrido que incluye renta variable, renta fija y ETFs regulados sobre criptoactivos.
    \item \textbf{Modelado del entorno de simulación} bajo el estándar Gymnasium, formalizando el problema como un Proceso de Decisión de Markov con costes de transacción y \textit{slippage} explícitos.
    \item \textbf{Entrenamiento del agente DRL} mediante el algoritmo \textit{Proximal Policy Optimization} (PPO), con una función de recompensa multiobjetivo que penaliza el \textit{Maximum Drawdown} y el \textit{turnover}.
    \item \textbf{Validación temporal estricta} mediante esquemas \textit{walk-forward} y \textit{expanding window}, junto con análisis de sensibilidad y por régimen de volatilidad.
    \item \textbf{Industrialización mediante una arquitectura de microservicios} que garantiza reproducibilidad, trazabilidad y aislamiento entre fases experimentales.
\end{enumerate}

El detalle técnico de cada fase —diseño del entorno, formulación de la recompensa, perfiles de riesgo, protocolo de validación y arquitectura del sistema— se desarrolla en el capítulo \ref{cap:materiales_y_metodos}, dedicado específicamente a Materiales y Métodos.

\subsection{Herramientas y Tecnologías}
El desarrollo del proyecto se apoya en el siguiente \textit{stack} tecnológico:

\begin{table}[h]
\centering
\begin{tabular}{|l|p{8cm}|}
\hline
\textbf{Categoría} & \textbf{Tecnología} \\ \hline
Lenguaje de Programación & Python 3.10+ \\ \hline
Framework de DRL & Stable-Baselines3 / Gymnasium \\ \hline
Análisis y Cálculo & Pandas / NumPy / PyTorch \\ \hline
Ingesta de Datos & Yahoo Finance API / \texttt{curl\_cffi} \\ \hline
Infraestructura y MLOps & FastAPI (Docker previsto como trabajo futuro) \\ \hline
\end{tabular}
\caption{Stack tecnológico del proyecto.}
\label{tab:herramientas}
\end{table}

\section{Planificación y Plan de Investigación}

El desarrollo del proyecto se estructurará en cuatro fases mensuales, integrando hitos de validación temprana para asegurar la robustez del sistema:

\begin{itemize}
    \item \textbf{Mes 1: Adquisición de Datos y Análisis Exploratorio.} 
    Selección del universo de activos (ETFs/ETPs). Implementación de la extracción automatizada mediante la API de \textit{Yahoo Finance} y preprocesamiento de series temporales (limpieza de \textit{outliers}, normalización e ingeniería de características/indicadores técnicos).
    
    \item \textbf{Mes 2: Construcción y Validación del Entorno Gymnasium.} 
    Desarrollo del entorno de simulación personalizado bajo el estándar \textit{Gymnasium}. \textbf{Hito Crítico:} Validación del entorno mediante \textit{baselines} simples (agente aleatorio y cartera equiponderada). Esta fase asegura que las métricas de recompensa, costes de transacción y lógica de rebalanceo funcionan correctamente antes de la fase de aprendizaje profundo.
    
    \item \textbf{Mes 3: Implementación y Entrenamiento del Agente DRL.} 
    Configuración y entrenamiento del algoritmo \textit{Proximal Policy Optimization} (PPO) utilizando la librería \textbf{Stable-Baselines3}. Esta elección permite centrar el esfuerzo técnico en el ajuste de hiperparámetros y el refinamiento de la función de recompensa personalizada, garantizando la estabilidad del entrenamiento.
    
    \item \textbf{Mes 4: Evaluación Comparativa, Análisis de Régimen y Redacción.} 
    Ejecución del protocolo de \textit{backtesting out-of-sample}. Comparativa detallada contra los \textit{baselines} financieros (60/40, Markowitz). Análisis de la capacidad de adaptación del agente ante cambios de régimen de mercado y redacción de las conclusiones finales de la memoria.
\end{itemize}

A continuación se presenta la planificación temporal del proyecto, desglosada por las fases metodológicas descritas y sus hitos críticos asociados.

\begin{figure}[h]
\centering
\begin{ganttchart}[
    expand chart=\textwidth,
    vgrid,
    hgrid,
    title/.style={fill=blue!10, draw=black, font=\bfseries},
    bar/.style={fill=blue!30, draw=black},
    bar height=0.6,
    progress label text={},
    group right shift=0,
    group top shift=0.7,
    group height=.3,
    milestone/.style={fill=red, draw=red, shape=diamond},
    y unit chart=0.6cm, % Ajustado para reducir la altura total del objeto
    title height=1
    ]{1}{16}
    
    \gantttitle{Plan de Investigación - 4 Meses}{16} \\
    \gantttitle{Mes 1}{4} \gantttitle{Mes 2}{4} \gantttitle{Mes 3}{4} \gantttitle{Mes 4}{4} \\
    \gantttitle{S1}{1} \gantttitle{S2}{1} \gantttitle{S3}{1} \gantttitle{S4}{1} 
    \gantttitle{S5}{1} \gantttitle{S6}{1} \gantttitle{S7}{1} \gantttitle{S8}{1} 
    \gantttitle{S9}{1} \gantttitle{S10}{1} \gantttitle{S11}{1} \gantttitle{S12}{1} 
    \gantttitle{S13}{1} \gantttitle{S14}{1} \gantttitle{S15}{1} \gantttitle{S16}{1} \\

    % Fase 1 - Asignamos nombre 'f1'
    \ganttbar[name=f1]{F1: Ingesta y Feature Engineering}{1}{4} \\
    \ganttmilestone{Dataset Validado}{4} \\
    
    % Fase 2 - Asignamos nombre 'f2'
    \ganttbar[name=f2]{F2: Entorno Gymnasium}{5}{8} \\
    \ganttmilestone{Hito: Validación Baselines}{8} \\
    
    % Fase 3 - Asignamos nombre 'f3'
    \ganttbar[name=f3]{F3: Entrenamiento DRL PPO}{9}{12} \\
    \ganttmilestone{Modelo Optimizado}{12} \\
    
    % Fase 4 - Asignamos nombre 'f4'
    \ganttbar[name=f4]{F4: Evaluación y Backtesting}{13}{15} \\
    \ganttbar{F5: Redacción Memoria}{13}{16} \\
    \ganttmilestone{Entrega Final}{16}
    
    % Enlaces corregidos usando los nombres asignados arriba
    \ganttlink{f1}{f2}
    \ganttlink{f2}{f3}
    \ganttlink{f3}{f4}
\end{ganttchart}
\caption{Cronograma detallado del proyecto por semanas.}
\label{fig:gantt_tfm}
\end{figure}

\section{Vinculación con los Objetivos de Desarrollo Sostenible (ODS)}
El presente proyecto se alinea principalmente con los siguientes objetivos de la Agenda 2030:
\begin{itemize}
    \item \textbf{ODS 8: Trabajo decente y crecimiento económico.} Al desarrollar herramientas que mejoran la eficiencia de los mercados financieros y la gestión del riesgo, se contribuye a la estabilidad del sistema económico y al crecimiento de soluciones Fintech innovadoras.
    \item \textbf{ODS 9: Industria, innovación e infraestructura.} La investigación en Inteligencia Artificial aplicada a las finanzas fomenta la innovación tecnológica y el desarrollo de infraestructuras financieras más resilientes y automatizadas.
    \item \textbf{ODS 12: Producción y consumo responsables.} La adopción de marcos de decisión cuantitativos y reproducibles favorece una gestión de carteras basada en evidencia y trazable, reduciendo la dependencia de decisiones discrecionales y contribuyendo a un uso más eficiente del capital invertido.
\end{itemize}


\chapter{Estado del Arte}
\label{cap:estado_del_arte}

Este capítulo ofrece una revisión crítica de la literatura científica y técnica que fundamenta la presente investigación. El objetivo es analizar la transición desde los modelos de optimización estática hacia los agentes de aprendizaje por refuerzo profundo (DRL) aplicados a la gestión de carteras, prestando especial atención a la integración de los activos digitales como nueva clase de inversión. La revisión se estructura en tres ejes temáticos: la evolución de los marcos de optimización de carteras, los avances recientes en DRL aplicado a finanzas y la institucionalización de los criptoactivos. El capítulo cierra con una síntesis crítica que delimita las limitaciones persistentes en la literatura y enuncia, de forma sintética, la novedad que aporta este trabajo.

\section{Evolución de la Gestión de Carteras: Del Paradigma de Markowitz a la Inteligencia Artificial}

La gestión cuantitativa de carteras ha transitado desde una disciplina basada en la intuición hacia un campo dominado por la optimización matemática y, más recientemente, por modelos dinámicos basados en aprendizaje automático. Esta evolución se enmarca en dos grandes etapas: una clásica, centrada en la eficiencia estadística estática, y una moderna, orientada al aprendizaje adaptativo bajo no estacionariedad.

\subsection{La Teoría Moderna de Carteras (MPT) y su Legado}
El punto de partida de la gestión cuantitativa moderna es el trabajo de Markowitz \cite{markowitz1952}. Su aportación fundamental reside en mostrar que el riesgo de un activo no debe medirse de forma aislada, sino a través de su contribución a la varianza total de la cartera. La MPT formaliza el problema de selección de activos como una optimización cuadrática en la que, dado un retorno esperado objetivo, se minimiza la varianza de la cartera.

Formalmente, la cartera óptima en sentido de Markowitz es la solución del problema:
\begin{equation}
\min_{\mathbf{w} \in \mathbb{R}^{n}} \; \mathbf{w}^{\top} \boldsymbol{\Sigma}\, \mathbf{w}
\quad \text{sujeto a} \quad
\mathbf{w}^{\top} \boldsymbol{\mu} = R^{*},
\quad \mathbf{w}^{\top} \mathbf{1} = 1,
\quad \mathbf{w} \geq \mathbf{0},
\label{eq:markowitz}
\end{equation}
donde $\mathbf{w} \in \mathbb{R}^{n}$ es el vector de pesos sobre $n$ activos, $\boldsymbol{\Sigma}$ es la matriz de covarianzas de los retornos, $\boldsymbol{\mu}$ el vector de retornos esperados, $R^{*}$ el retorno objetivo y $\mathbf{1}$ el vector unitario. La restricción $\mathbf{w}^{\top}\mathbf{1}=1$ garantiza que la cartera está plenamente invertida, mientras que $\mathbf{w}\geq\mathbf{0}$ corresponde a un escenario \textit{long-only}, habitual en la gestión institucional regulada.

Pese a su valor fundacional, la aplicación práctica de la MPT presenta limitaciones ampliamente documentadas en la literatura sobre eficiencia y racionalidad de los mercados \cite{michaud1989}:

\begin{itemize}
    \item \textbf{Sensibilidad a los parámetros estimados.} La solución óptima depende fuertemente de las estimaciones de $\boldsymbol{\mu}$ y $\boldsymbol{\Sigma}$, que se calculan sobre ventanas históricas finitas. Pequeñas variaciones en estos parámetros producen reasignaciones drásticas de pesos, lo que cuestiona la estabilidad práctica del modelo.
    \item \textbf{Supuesto implícito de estacionariedad.} La MPT trata las correlaciones como constantes a lo largo del tiempo. La evidencia empírica muestra, sin embargo, que en periodos de tensión las correlaciones tienden a aumentar, reduciendo el beneficio de la diversificación previsto por el modelo.
    \item \textbf{Distribución gaussiana de los retornos.} Numerosos estudios documentan que los retornos financieros presentan colas pesadas y exceso de curtosis, propiedades que la suposición de normalidad del marco clásico no captura adecuadamente.
\end{itemize}

\subsection{Hacia una Optimización Secuencial: Limitaciones del Marco Estático}

Las limitaciones anteriores se acentúan cuando la cartera incorpora activos cuyo comportamiento estadístico se aleja de los supuestos clásicos, como ocurre con los criptoactivos. La literatura ha respondido extendiendo el marco original con diversas variantes (Black-Litterman, optimización robusta, modelos de mínima varianza con restricciones de norma, etc.), pero todas comparten la naturaleza estática del enfoque: cada decisión se toma resolviendo un problema de optimización de un único periodo, asumiendo parámetros estimados a partir de una ventana retrospectiva.

Frente a este paradigma estático, una línea creciente de literatura plantea la gestión de carteras como un problema secuencial en el que la decisión actual condiciona el espacio de estados futuro. Bajo esta visión, las comisiones, el \textit{slippage} y el \textit{turnover} dejan de ser fricciones externas y pasan a formar parte intrínseca de la función objetivo. Este replanteamiento abre la puerta a marcos basados en procesos de decisión de Markov y, en particular, al aprendizaje por refuerzo profundo. La tabla \ref{tab:comparativa_paradigmas} resume el contraste entre el paradigma estático clásico, los enfoques basados en DRL y la contribución específica de los ecosistemas de estandarización tipo FinRL.

\begin{table}[h]
\centering
\renewcommand{\arraystretch}{1.25}
\begin{tabular}{|p{3.0cm}|p{3.7cm}|p{3.7cm}|p{3.7cm}|}
\hline
\textbf{Dimensión} & \textbf{MPT clásica (Markowitz)} & \textbf{DRL aplicado a carteras} & \textbf{Ecosistema FinRL} \\ \hline
Naturaleza de la decisión & Estática, un único periodo & Secuencial, política condicionada al estado & Hereda la naturaleza secuencial del DRL \\ \hline
Tratamiento de fricciones & Externo al modelo & Integrado en la recompensa & Soporte parcial vía wrappers genéricos \\ \hline
Supuesto de estacionariedad & Sí (correlaciones constantes) & No (la política puede adaptarse a cambios de régimen) & No impone supuestos adicionales \\ \hline
Adecuación a colas pesadas & Limitada (asume normalidad) & Adecuada (no requiere supuestos distribucionales) & Heredada del algoritmo DRL empleado \\ \hline
Reproducibilidad & Alta (forma cerrada) & Baja sin estandarización; sensible a semillas y datos & Aporte central: estandariza ingesta, entornos y \textit{backtesting} \\ \hline
Ámbito de aplicación habitual & Carteras institucionales tradicionales & Estudios mayoritariamente sobre acciones o cripto puro & Bursátil, trading de un único activo o universos pequeños \\ \hline
Fricción operativa con criptoactivos regulados (ETFs) & No abordada explícitamente & Posible pero requiere diseño \textit{ad-hoc} & No incluida en los entornos de referencia \\ \hline
\end{tabular}
\caption{Contraste sintético entre el paradigma estático de Markowitz, los enfoques basados en DRL y los ecosistemas de estandarización tipo FinRL.}
\label{tab:comparativa_paradigmas}
\end{table}

Como pone de manifiesto la tabla \ref{tab:comparativa_paradigmas}, ninguno de los tres marcos resuelve por sí solo el problema planteado en este trabajo: la MPT clásica no captura la dinámica de los criptoactivos, los enfoques de DRL aportan flexibilidad pero rara vez se aplican a universos híbridos, y los ecosistemas tipo FinRL aportan reproducibilidad pero sin un soporte específico para ETFs sobre criptoactivos ni para funciones de recompensa multiobjetivo. Las dos secciones siguientes profundizan, primero, en los avances específicos del DRL aplicado a carteras y, después, en la integración de los criptoactivos como nueva clase de activo.

\section{Deep Reinforcement Learning (DRL) en la Toma de Decisiones Financieras}

El \textit{Deep Reinforcement Learning} aplicado a la gestión de carteras supone un cambio de paradigma respecto a los enfoques estáticos: la decisión deja de obtenerse como solución de un problema de optimización de un único periodo y pasa a aprenderse como política condicionada al estado del entorno. La literatura modela este escenario como un Proceso de Decisión de Markov (MDP), donde un agente interactúa con el mercado, observa un estado $\mathbf{s}_t$, ejecuta una acción $\mathbf{a}_t$ y recibe una recompensa $r_t$, con el objetivo de maximizar el retorno acumulado esperado.

\subsection{El Algoritmo PPO como Estándar de Estabilidad}

Dentro de la familia de algoritmos de gradiente de política, \textit{Proximal Policy Optimization} (PPO) \cite{schulman2017ppo} se ha consolidado como una elección de referencia debido a su estabilidad y eficiencia muestral. Una de las dificultades clásicas de los métodos de gradiente de política es que actualizaciones excesivamente amplias pueden conducir a degradaciones bruscas en el desempeño, especialmente en entornos con baja relación señal-ruido como los financieros. PPO mitiga este problema introduciendo un objetivo recortado (\textit{Clipped Surrogate Objective}) que limita la magnitud del cambio entre la política actualizada y la previa.

Formalmente, sea $\pi_{\theta}$ la política parametrizada y $\hat{A}_t$ una estimación de la ventaja en el instante $t$. El objetivo recortado se define como:
\begin{equation}
L^{\mathrm{CLIP}}(\theta) = \hat{\mathbb{E}}_t \!\left[ \min\!\Big( r_t(\theta)\,\hat{A}_t,\; \mathrm{clip}\!\big(r_t(\theta),\, 1-\epsilon,\, 1+\epsilon\big)\,\hat{A}_t \Big) \right],
\label{eq:ppo-clip}
\end{equation}
donde
\begin{equation}
r_t(\theta) = \frac{\pi_{\theta}(\mathbf{a}_t \mid \mathbf{s}_t)}{\pi_{\theta_{\mathrm{old}}}(\mathbf{a}_t \mid \mathbf{s}_t)}
\end{equation}
es el cociente de probabilidades entre la política nueva y la antigua, y $\epsilon$ es un hiperparámetro que acota el cambio relativo permitido en cada actualización. El operador $\min$ aplicado al producto $r_t(\theta)\,\hat{A}_t$ y a su versión recortada penaliza los movimientos que se alejan del entorno de confianza implícito definido por $[1-\epsilon, 1+\epsilon]$, evitando actualizaciones inestables sin necesidad de imponer restricciones explícitas sobre la divergencia de Kullback-Leibler como en el caso de TRPO.

Esta estabilidad ha favorecido la adopción de PPO en múltiples estudios sobre trading algorítmico, en los que la combinación de alta varianza de la señal y horizontes largos de inversión hace especialmente sensible al método de optimización empleado.

\subsubsection{Comparación con algoritmos alternativos}
\label{sec:ppo_vs_alternativas}

La elección de PPO no es la única opción posible dentro del DRL, y su justificación frente a otros algoritmos requiere examinarse a la luz de los requisitos específicos del problema: espacio de acción continuo de alta dimensionalidad, baja relación señal-ruido y necesidad de estabilidad en el entrenamiento sobre series financieras. La tabla \ref{tab:ppo_vs_alternativas} sintetiza el contraste con las alternativas más habituales.

\begin{table}[h]
\centering
\renewcommand{\arraystretch}{1.25}
\begin{tabular}{|p{2.6cm}|p{6.0cm}|p{5.2cm}|}
\hline
\textbf{Algoritmo} & \textbf{Características relevantes} & \textbf{Limitaciones para gestión de carteras} \\ \hline
DQN \cite{schulman2017ppo} & Aprendizaje basado en función de valor; espacio de acción discreto. & Inadecuado para asignación continua de pesos en $[0,1]^{n}$; requeriría discretizar el simplex con pérdida de granularidad. \\ \hline
A2C & Actor-Crítico síncrono; un único paso de actualización por lote. & Menor eficiencia muestral que PPO al no reutilizar trayectorias mediante \textit{minibatch SGD} multi-época; mayor varianza de gradiente. \\ \hline
DDPG & Política determinista para acciones continuas; aprendizaje \textit{off-policy}. & Sensibilidad elevada a hiperparámetros y a la inicialización; rendimiento frágil en entornos estocásticos como los financieros. \\ \hline
SAC & \textit{Off-policy}, actor-crítico estocástico, regularización por entropía. & Buen rendimiento empírico, pero coste computacional y de memoria mayor (\textit{replay buffer}, doble crítico) y sintonización más exigente. \\ \hline
TRPO \cite{schulman2017ppo} & Restricción explícita de divergencia KL para garantizar pasos seguros. & Implementación notablemente más compleja (segunda derivada, conjugado-gradiente) sin mejora consistente respecto a PPO. \\ \hline
PPO \cite{schulman2017ppo} & Objetivo recortado, \textit{on-policy}, actor-crítico estocástico, eficiencia muestral mediante \textit{minibatch SGD} multi-época. & Requiere recolección de trayectorias frescas en cada actualización; el coste se compensa por la estabilidad y la simplicidad de la implementación. \\ \hline
\end{tabular}
\caption{Comparativa de PPO frente a alternativas de DRL desde la perspectiva de la gestión secuencial de carteras.}
\label{tab:ppo_vs_alternativas}
\end{table}

A partir de la comparativa anterior, la elección de PPO en el presente trabajo se justifica por tres motivos. En primer lugar, soporta de forma nativa espacios de acción continuos sin necesidad de discretizar el simplex de pesos, lo que evita una pérdida indeseada de granularidad en la asignación. En segundo lugar, su mecanismo de objetivo recortado proporciona estabilidad de entrenamiento sin la complejidad de implementación de TRPO ni la fragilidad de DDPG. Por último, en términos de coste de implementación y reproducibilidad, PPO ha sido objeto de un esfuerzo de estandarización particularmente intenso —en especial a través de librerías como Stable-Baselines3—, lo que reduce la probabilidad de error en el código de entrenamiento y facilita la replicabilidad de los experimentos por terceros.

\subsection{Arquitectura Actor-Crítico}

La mayoría de implementaciones modernas de PPO se apoyan en una arquitectura \textit{Actor-Critic}, en la que dos redes neuronales colaboran:
\begin{itemize}
    \item El \textit{actor} aproxima la política $\pi_{\theta}(\mathbf{a}\mid\mathbf{s})$ y propone la acción a ejecutar (en gestión de carteras, el vector de pesos $\mathbf{w}_t$).
    \item El \textit{crítico} estima la función de valor $V_{\phi}(\mathbf{s})$, lo que permite calcular ventajas $\hat{A}_t$ con menor varianza que un estimador puramente Monte Carlo.
\end{itemize}
Esta separación es especialmente útil en problemas con espacios de acción continuos, como la asignación de pesos de cartera en $[0,1]^{n}$, donde aproximaciones discretas (por ejemplo, basadas en DQN) resultan poco adecuadas.

\subsection{Industrialización del DRL Financiero: el Ecosistema FinRL}

Una dificultad recurrente en la literatura de DRL financiero es la reproducibilidad: muchos resultados publicados son difíciles de replicar debido a la variabilidad de fuentes de datos, preprocesamientos y parametrizaciones del entorno. Para mitigar esta brecha, los trabajos en torno a \textbf{FinRL} \cite{liu2018finrl,liu2021finrl} proponen una infraestructura jerárquica que estandariza la ingesta de datos, el preprocesamiento y el \textit{backtesting} sobre el estándar \textbf{Gymnasium}.

Dicha estandarización tiene un valor doble. Por un lado, reduce la fricción de implementación al ofrecer entornos preconfigurados que cumplen la interfaz Gymnasium, permitiendo encadenar algoritmos de la familia PPO/A2C/DDPG sin reescribir el simulador. Por otro lado, fija un marco común para reportar métricas de evaluación, lo que facilita la comparación entre trabajos.

A pesar de su valor como infraestructura de referencia, conviene delimitar con precisión qué aporta FinRL y qué queda fuera de su alcance, dado que esa frontera condiciona directamente el diseño del presente trabajo. La tabla \ref{tab:finrl_vs_tfm} sintetiza esa delimitación.

\begin{table}[h]
\centering
\renewcommand{\arraystretch}{1.25}
\begin{tabular}{|p{4.0cm}|p{5.0cm}|p{5.0cm}|}
\hline
\textbf{Componente} & \textbf{Aportación de FinRL} & \textbf{A desarrollar en este trabajo} \\ \hline
Ingesta de datos & Conectores estandarizados a varias fuentes (Yahoo, Alpaca, etc.) & Ingesta sobre Yahoo en entornos restrictivos (\texttt{curl\_cffi}) y selección de universo orientada a una cartera híbrida \\ \hline
Entornos de simulación & Entornos genéricos para \textit{trading} de acciones, criptomonedas y carteras pequeñas & Entorno específico para cartera híbrida con ETFs regulados sobre criptoactivos (IBIT, ETHA), restricción \textit{long-only} y suma a uno \\ \hline
Función de recompensa & Recompensas estándar basadas en retorno o ratio de Sharpe sobre ventana fija & Recompensa multiobjetivo que combina Sharpe rolling, penalización por \textit{Maximum Drawdown} y por \textit{turnover}, parametrizada por perfiles de riesgo \\ \hline
Modelado de fricciones & Soporte de comisiones lineales en algunos entornos & Modelado explícito de comisiones y \textit{slippage}, con reflejo directo en la dinámica del entorno \\ \hline
Protocolo de validación & Particiones \textit{train/test} simples sobre periodo fijo & Validación temporal estricta con esquemas \textit{walk-forward} y \textit{expanding window}, análisis de sensibilidad y por régimen de volatilidad \\ \hline
Trazabilidad experimental & Registro mínimo de experimentos en notebooks & Persistencia en base de datos relacional con metadatos completos por modelo entrenado y reproducción exacta de simulaciones \\ \hline
Despliegue & No abordado de forma específica & Capa de servicios FastAPI, autenticación, roles y arquitectura de microservicios para uso operativo \\ \hline
\end{tabular}
\caption{Delimitación entre lo que aporta FinRL como infraestructura y lo que se desarrolla específicamente en este trabajo.}
\label{tab:finrl_vs_tfm}
\end{table}

En síntesis, FinRL aporta el \textit{andamiaje} —ingesta y \textit{wrappers} sobre Gymnasium— pero deja abiertas precisamente las cuestiones que resultan críticas para la gestión institucional de carteras híbridas: el modelado fino de fricciones, la inclusión de instrumentos regulados sobre criptoactivos y el diseño de funciones de recompensa que penalicen explícitamente \textit{drawdown} y \textit{turnover}. Trabajos como \cite{jiang2017drl}, que aplican DRL al problema concreto de gestión de cartera, ilustran el potencial del enfoque secuencial pero se centran en universos de criptomonedas \textit{spot} sin integrar instrumentos regulados ni penalizaciones explícitas de riesgo y rotación. La sección siguiente examina precisamente las propiedades estadísticas de los criptoactivos y el papel del IBIT como instrumento regulado, cuya integración constituye uno de los focos del presente trabajo.

\section{Criptoactivos e Institucionalización: el Impacto del IBIT}

La integración de criptoactivos en carteras institucionales ha trascendido el ámbito puramente especulativo y se ha consolidado como un componente con potencial de diversificación estructural. Como se argumenta en \cite{hougan2021crypto}, el Bitcoin presenta propiedades estadísticas que lo distinguen tanto de la renta variable tradicional como de los activos refugio clásicos (oro o bonos), aunque sus efectos en términos de diversificación dependen fuertemente del régimen de mercado.

\subsection{El IBIT como Catalizador de Eficiencia y Liquidez}

El lanzamiento del \textit{iShares Bitcoin Trust} (IBIT) constituye un punto de inflexión en la microestructura del mercado de criptoactivos. Al tratarse de un ETF \textit{spot} regulado, el IBIT permite acceder a la exposición a Bitcoin a través de los canales habituales de negociación bursátil, eliminando una parte significativa de las fricciones de custodia y liquidez asociadas a los mercados \textit{offshore}.

Desde la perspectiva de la teoría de carteras, la inclusión de un activo con un perfil de retorno-riesgo diferenciado podría contribuir a desplazar la frontera eficiente, mejorando el retorno por unidad de riesgo. No obstante, este beneficio es contingente al régimen de mercado: la literatura previa documenta episodios en los que el Bitcoin ha mostrado un acoplamiento elevado con índices bursátiles tradicionales como el S\&P 500, especialmente durante eventos de tensión financiera. \cite{conlon2020safehaven} muestran, sobre el periodo de mercado bajista asociado a la COVID-19, que el Bitcoin no se comportó como activo refugio sino que su correlación con la renta variable aumentó en los momentos de mayor estrés. Esta dependencia condicional sugiere que la diversificación obtenida mediante criptoactivos no es estable a lo largo del tiempo, lo que motiva el uso de marcos de gestión dinámica capaces de adaptar la exposición al régimen vigente.

\subsection{Propiedades Estadísticas Relevantes para la Gestión Algorítmica}

Los retornos del IBIT heredan del Bitcoin un conjunto de propiedades estadísticas relevantes desde el punto de vista del modelado:

\begin{itemize}
    \item \textbf{\textit{Volatility clustering}.} Tal y como se documenta en \cite{katsiampa2017}, la volatilidad del Bitcoin presenta una marcada estructura de \textit{clustering}, ajustándose adecuadamente a modelos de la familia GARCH. Bajo este patrón, periodos de alta agitación alternan con fases comparativamente más estables, en lugar de distribuirse uniformemente en el tiempo.
    \item \textbf{Asimetría y colas pesadas.} Aunque la serie histórica del IBIT es todavía corta (cotiza desde enero de 2024) y no permite por sí sola caracterizaciones distribucionales robustas, los retornos del Bitcoin subyacente presentan, según la literatura, asimetría y colas significativamente más pesadas que las de una distribución normal, propiedades que cabe esperar igualmente en el ETF y que implican una mayor frecuencia de eventos extremos que la prevista por los modelos gaussianos clásicos.
    \item \textbf{Dependencia condicional con el mercado tradicional.} Como se ha indicado, las correlaciones con renta variable y renta fija varían con el régimen, lo que penaliza a los modelos estáticos que asumen una estructura de covarianza constante.
\end{itemize}

Estas características, combinadas con la naturaleza no estacionaria de los retornos, dificultan la aplicación directa de los marcos de optimización de un único periodo y refuerzan el interés de aproximaciones secuenciales basadas en aprendizaje por refuerzo profundo.

\section{Síntesis Crítica y Hueco de Investigación}
\label{sec:sintesis_critica}

La revisión de los apartados anteriores permite extraer tres observaciones principales:

\begin{enumerate}
    \item Los marcos clásicos de optimización de carteras, basados en el legado de Markowitz, asumen estacionariedad, normalidad y ausencia de fricciones, supuestos que se ven especialmente comprometidos cuando la cartera incorpora criptoactivos.
    \item Los algoritmos de DRL —y en particular PPO— ofrecen un marco maduro para abordar el problema como una decisión secuencial bajo incertidumbre. Sin embargo, gran parte de la literatura aplicada se centra en universos puramente bursátiles o en criptomonedas aisladas, y emplea funciones de recompensa relativamente simples (típicamente, retornos logarítmicos o ratios de Sharpe sobre ventana fija).
    \item Ecosistemas como FinRL han contribuido decisivamente a estandarizar la ingesta y la evaluación, pero no abordan de forma específica el problema de gestionar carteras híbridas que combinen activos tradicionales con ETFs sobre criptoactivos, ni el de reflejar con fidelidad las fricciones operativas asociadas a tales instrumentos.
\end{enumerate}

A partir de estas observaciones, el hueco de investigación que motiva el presente trabajo puede formularse como sigue: \textit{no existe en la literatura revisada un marco integrado que combine (i) un agente DRL basado en PPO, (ii) una función de recompensa multiobjetivo que penalice explícitamente \textnormal{drawdown} y \textnormal{turnover}, (iii) un universo híbrido que incluya ETFs regulados sobre Bitcoin y Ethereum, y (iv) un protocolo de validación temporal estricto sobre datos out-of-sample}. Cubrir este vacío constituye la motivación central del TFM y se concreta en las contribuciones que se sintetizan a continuación.

\section{Novedad de este Trabajo}
\label{sec:novedad}

A partir del hueco identificado en la sección \ref{sec:sintesis_critica}, este Trabajo Final de Máster se posiciona en la intersección entre el aprendizaje por refuerzo profundo, la gestión institucional de carteras y la incorporación de activos digitales como nueva clase de inversión. Sus aportaciones, que se desarrollan en detalle en los capítulos posteriores, pueden sintetizarse en los siguientes ejes:

\begin{itemize}
    \item Diseño de un agente PPO sobre un MDP específicamente adaptado a una cartera híbrida que combina renta variable tradicional, renta fija y ETFs regulados sobre criptoactivos (IBIT y ETHA).
    \item Definición de una función de recompensa multiobjetivo que integra retorno ajustado por riesgo, penalización por \textit{drawdown} y penalización por \textit{turnover}, parametrizada mediante perfiles de riesgo seleccionables.
    \item Aplicación de un protocolo de validación temporal robusto, basado en \textit{walk-forward}, \textit{expanding window} y análisis de sensibilidad de los hiperparámetros de la recompensa, así como un análisis del rendimiento por régimen de volatilidad.
    \item Implementación del sistema como una arquitectura de microservicios con persistencia, autenticación y trazabilidad de cada modelo entrenado, lo que permite la reproducción exacta de los experimentos y la simulación personalizada por parte de usuarios finales.
\end{itemize}

El detalle de cómo se concretan estas aportaciones —el modelado del entorno, la metodología de validación, la arquitectura del sistema y los experimentos realizados— se aborda en los capítulos siguientes de la memoria.

\chapter{Materiales y Métodos}
\label{cap:materiales_y_metodos}

Este capítulo describe los elementos metodológicos del trabajo: el marco general de investigación adoptado, la arquitectura del sistema implementado y las decisiones de ingeniería que sostienen el ciclo experimental. Las secciones que siguen presentan, primero, el encaje del trabajo en el marco de la \textit{Design Science Research}; a continuación, la arquitectura de microservicios diseñada para garantizar reproducibilidad y resiliencia; y, finalmente, las decisiones específicas relativas a la ingesta de datos en entornos restrictivos. El detalle del entorno DRL, la función de recompensa y el protocolo de validación se desarrolla en los capítulos posteriores dedicados a la implementación y a los resultados.

\section{Marco metodológico: Design Science Research}
\label{sec:dsr}

El presente trabajo se encuadra en el marco de la \textit{Design Science Research} (DSR) \cite{oates2006}, una metodología orientada a la creación y evaluación de artefactos tecnológicos como medio para resolver problemas prácticos relevantes. El \textit{artefacto} de esta investigación es un sistema integral de optimización de carteras híbridas basado en aprendizaje por refuerzo profundo, compuesto por un pipeline de datos, un entorno de simulación, un agente PPO entrenado y una capa de servicios para evaluación y simulación personalizada.

La adopción de DSR resulta adecuada porque el trabajo no se limita a verificar una hipótesis científica aislada, sino que persigue la construcción de un sistema completo cuya evaluación se realiza tanto en términos cuantitativos (métricas financieras de riesgo y rendimiento) como cualitativos (reproducibilidad, mantenibilidad y trazabilidad operativa). El ciclo de diseño se ha articulado en iteraciones sucesivas en las que cada fase aporta un componente del artefacto y un protocolo de evaluación específico, en línea con el esquema clásico propuesto en la literatura de DSR.

\section{Arquitectura del Sistema: Enfoque MLOps}
\label{sec:arquitectura}

La eficacia de un agente DRL depende, en la práctica, de una infraestructura de soporte capaz de garantizar reproducibilidad, trazabilidad y aislamiento entre fases experimentales. \cite{sculley2015hidden} señalan que la deuda técnica acumulada en el preprocesamiento, en el versionado de modelos y en la gestión de datos suele ser determinante en la capacidad de un sistema de aprendizaje automático para sostener resultados consistentes a lo largo del tiempo, hasta el punto de superar en magnitud a la complejidad estrictamente algorítmica. A partir de estas consideraciones, este trabajo adopta un enfoque de \textit{MLOps} en el que la arquitectura se descompone en microservicios coordinados, con persistencia explícita de datos, modelos y experimentos.

\subsection{Resiliencia mediante Microservicios y Contenedores}

La arquitectura se organiza en torno a una capa de servicio basada en \textbf{FastAPI}, sobre la que se exponen las distintas fases del pipeline (selección del universo, preparación de datos, entrenamiento del agente, validación temporal y simulación personalizada). Esta separación modular permite independizar el ciclo de inferencia del agente respecto a la ingesta y al preprocesamiento, así como escalar de forma diferenciada cada componente. El uso de contenedores con \textbf{Docker} se prevé para garantizar la paridad de entornos entre fase de entrenamiento y despliegue, evitando inconsistencias entre dependencias del ecosistema Gymnasium y de las librerías de \textit{Deep Learning}.

La capa de persistencia está soportada por una base de datos relacional gestionada mediante un ORM, en la que se almacenan los universos de activos seleccionados por el screener, los modelos entrenados con su configuración de recompensa, y el historial de simulaciones realizadas por los usuarios. Esta capa, junto con la persistencia en disco de los CSVs de \textit{features} y los modelos serializados, constituye la base de la trazabilidad experimental del sistema.

\subsection{Ingesta de Datos en Entornos Restrictivos}

La ingesta de datos se ha diseñado teniendo en cuenta restricciones habituales en entornos corporativos, tales como la presencia de \textit{proxies} SSL y la posibilidad de bloqueos a las cabeceras estándar utilizadas por las librerías Python más comunes. Para mitigar estas dificultades, se ha integrado la librería \texttt{curl\_cffi}, que permite emular la huella TLS de un navegador Chrome moderno y, con ello, acceder a fuentes públicas como Yahoo Finance sin comprometer las políticas de seguridad de la red.

Esta solución se ha encapsulado tras una abstracción de fuente de datos (\textit{Strategy pattern}), de modo que el resto del pipeline trabaja contra una interfaz común y resulta independiente del proveedor concreto utilizado. De este modo, la arquitectura admite, sin modificar el resto del sistema, la sustitución futura de la fuente histórica por proveedores en tiempo real (por ejemplo, Alpaca, Polygon.io o Interactive Brokers), siempre que cumplan con la interfaz especificada.

\section{Universo de Activos}
\label{sec:universo}

El sistema opera sobre un universo híbrido de nueve activos seleccionados para reflejar tanto las clases tradicionales (renta variable, renta fija, sectores defensivos) como la nueva clase de activos digitales accesible mediante ETFs regulados. La tabla \ref{tab:universo} resume el universo y la naturaleza de cada inclusión.

\begin{table}[h]
\centering
\renewcommand{\arraystretch}{1.25}
\begin{tabular}{|l|l|p{6.5cm}|}
\hline
\textbf{Ticker} & \textbf{Clase} & \textbf{Función en la cartera} \\ \hline
IVV & Renta variable USA & Réplica del S\&P 500; proxy del mercado bursátil estadounidense. \\ \hline
BND & Renta fija agregada & ETF de Vanguard sobre el mercado total de bonos USA. \\ \hline
IBIT & Cripto regulada & ETF \textit{spot} de Bitcoin (iShares); foco principal del TFM. \\ \hline
ETHA & Cripto regulada & ETF \textit{spot} de Ethereum; segunda exposición cripto del universo. \\ \hline
MO & Defensivo (tabaco) & Activo defensivo de alta dividend yield. \\ \hline
JNJ & Defensivo (salud) & Defensivo de gran capitalización con dividendo estable. \\ \hline
SCU & Diversificación & Activo del sector financiero alternativo. \\ \hline
AWK & Defensivo (utilities) & Servicios públicos regulados. \\ \hline
CB & Defensivo (seguros) & Aseguradora global de gran capitalización. \\ \hline
\end{tabular}
\caption{Universo de activos del sistema.}
\label{tab:universo}
\end{table}

La selección busca representar una cartera realista de inversión institucional sobre la que evaluar la integración de criptoactivos: dispone de un núcleo amplio (IVV/BND), una exposición cripto diferenciada (IBIT y ETHA, dos clases con riesgo similar pero dinámica distinta) y una capa de activos defensivos sectoriales que aportan baja correlación con los activos de riesgo en periodos de tensión.

\paragraph{Limitación conocida.} La construcción actual de las correlaciones \textit{rolling} como variable de estado del entorno utiliza tres pares (IBIT--IVV, IBIT--BND, IVV--BND). ETHA, aunque sí participa como activo asignable de la cartera, no aparece en ninguno de los pares de correlación dinámica del estado, lo que constituye una limitación documentada en la implementación cuya extensión queda como trabajo futuro.

\section{Ingeniería de Características}
\label{sec:features}

A partir de las series de precios de cierre de los nueve activos del universo, se construye un conjunto de características diseñadas para que el agente disponga de información sobre la dinámica reciente, la volatilidad y los regímenes de mercado. Las categorías y ventanas se resumen en la tabla \ref{tab:features}.

\begin{table}[h]
\centering
\renewcommand{\arraystretch}{1.25}
\begin{tabular}{|p{3.0cm}|p{3.7cm}|p{6.3cm}|}
\hline
\textbf{Categoría} & \textbf{Indicadores} & \textbf{Ventanas y propósito} \\ \hline
Retornos & Log-retornos diarios & Magnitud de variación día a día. \\ \hline
Momentum & Retorno acumulado a 5d, 20d y 60d & Tendencia a corto, medio y largo plazo, en línea con los horizontes documentados en la literatura de \textit{factor momentum}. \\ \hline
Tendencia y momentum técnico & RSI, MACD diff & Sobrecompra/sobreventa y diferencia entre medias móviles. \\ \hline
Volatilidad & ATR, \textit{Bollinger Band Width}, desviación estándar \textit{rolling}, asimetría y curtosis & Régimen de volatilidad y forma de la distribución reciente. \\ \hline
Oscilador de canal & CCI & Distancia al canal típico de precio. \\ \hline
Correlaciones & Correlación \textit{rolling} entre pares (IBIT--IVV, IBIT--BND, IVV--BND) & Estructura dinámica de dependencia entre activos. \\ \hline
\end{tabular}
\caption{Conjunto de características que componen el espacio de observación del agente.}
\label{tab:features}
\end{table}

Antes de alimentar al agente, las \textit{features} se normalizan por activo y se truncan los \textit{outliers} para estabilizar el entrenamiento. El conjunto resultante se almacena en disco como \texttt{normalized\_features.csv} y constituye, junto con \texttt{original\_prices.csv}, la entrada estática del entorno de simulación.

\section{Entorno de Simulación: \texttt{PortfolioEnv}}
\label{sec:entorno}

El entorno de simulación se ha implementado siguiendo la interfaz Gymnasium estándar y formaliza el problema como un Proceso de Decisión de Markov (MDP) finito. Sus componentes son:

\begin{itemize}
    \item \textbf{Espacio de estados} $\mathcal{S}$. Vector de observación del día $t$ que concatena las \textit{features} normalizadas de los nueve activos para ese instante, los pesos actuales de la cartera $\mathbf{w}_{t-1}$ y la última recompensa observada $r_{t-1}$. Se trata de una observación aumentada que permite al agente disponer no solo del estado de mercado, sino también de su propia posición previa.
    \item \textbf{Espacio de acciones} $\mathcal{A}$. Vector continuo $\mathbf{a}_t \in [0,1]^{n}$ con $n=9$. La acción se interpreta como pesos no normalizados sobre el simplex de cartera: el entorno aplica internamente la normalización
    \begin{equation*}
        \mathbf{w}_t = \frac{\mathbf{a}_t}{\sum_i a_{t,i}},
    \end{equation*}
    garantizando $\mathbf{w}_t^{\top}\mathbf{1}=1$ y $\mathbf{w}_t \geq \mathbf{0}$ (cartera \textit{long-only}).
    \item \textbf{Dinámica del entorno}. Dado el vector de pesos $\mathbf{w}_t$, el entorno aplica el rebalanceo, descuenta comisiones de transacción proporcionales al \textit{turnover} y un \textit{slippage} parametrizable, y actualiza el valor de la cartera según los retornos del día. La transición hacia el siguiente estado consiste en avanzar un paso en la serie temporal de \textit{features}.
    \item \textbf{Recompensa}. Función multiobjetivo que combina retorno ajustado por riesgo y penalizaciones por riesgo realizado y rotación, descrita formalmente en la sección \ref{sec:recompensa}.
\end{itemize}

\section{Función de Recompensa y Perfiles de Riesgo}
\label{sec:recompensa}

\subsection{Formulación de la Recompensa}

La recompensa por paso del entorno se define como:
\begin{equation}
r_t = \mathrm{clip}\!\left(\,\mathrm{Sharpe}^{(20)}_t \;-\; \varphi \cdot \mathrm{MDD}_t \;-\; \gamma \cdot T_t\,,\; -1,\; 1\,\right),
\label{eq:reward}
\end{equation}
donde:
\begin{itemize}
    \item $\mathrm{Sharpe}^{(20)}_t$ es el ratio de Sharpe calculado sobre los retornos diarios de la cartera en una ventana móvil de 20 días, normalizado para acotarse en un rango similar a $[-1, 1]$.
    \item $\mathrm{MDD}_t \in [0, 1]$ es el \textit{Maximum Drawdown} actual respecto al máximo histórico del valor de cartera observado hasta el instante $t$.
    \item $T_t = \sum_i |w_{i,t} - w_{i,t-1}|$ es el \textit{turnover} del paso $t$, suma de los cambios absolutos en pesos.
    \item $\varphi$ y $\gamma$ son hiperparámetros que ponderan, respectivamente, la aversión al \textit{drawdown} y la aversión a la rotación. Su elección define el \textit{perfil de riesgo} del agente.
    \item El operador $\mathrm{clip}$ acota el reward en $[-1,1]$ para estabilizar el entrenamiento y evitar que valores extremos dominen el gradiente.
\end{itemize}

A diferencia de las recompensas basadas en log-retorno o ratio de Sharpe sobre ventana fija, esta formulación es \textit{multiobjetivo}: integra explícitamente la dimensión de riesgo (vía $\mathrm{MDD}$) y la dimensión de fricción operativa (vía $T_t$) en una única señal escalar, permitiendo al algoritmo de optimización tratarlas como objetivos contrapuestos a equilibrar.

\subsection{Perfiles de Riesgo}
\label{sec:perfiles}

Los hiperparámetros $\varphi$ y $\gamma$ se ofrecen al usuario administrador como cuatro perfiles preconfigurados (tabla \ref{tab:perfiles}), cada uno representando una filosofía de inversión distinta.

\begin{table}[h]
\centering
\renewcommand{\arraystretch}{1.25}
\begin{tabular}{|l|c|c|p{7.6cm}|}
\hline
\textbf{Perfil} & $\boldsymbol{\varphi}$ & $\boldsymbol{\gamma}$ & \textbf{Filosofía} \\ \hline
\texttt{balanced} & 0.02 & 0.010 & Equilibrio entre retorno y control de riesgo. \textit{Baseline} del catálogo. \\ \hline
\texttt{conservative} & 0.05 & 0.010 & Penaliza \textit{drawdowns} con factor 2.5x respecto a \texttt{balanced}; prioriza preservar capital. \\ \hline
\texttt{low\_turnover} & 0.02 & 0.020 & \textbf{Perfil principal del TFM.} Doble penalización por rotación que \texttt{balanced}: el agente solo opera cuando la mejora esperada en Sharpe compensa el coste de rotar. \\ \hline
\texttt{aggressive} & 0.01 & 0.005 & Mínimas penalizaciones; máxima libertad para buscar retorno. \\ \hline
\end{tabular}
\caption{Perfiles de riesgo predefinidos del catálogo.}
\label{tab:perfiles}
\end{table}

La elección de \texttt{low\_turnover} como perfil principal del TFM no es arbitraria: se justifica empíricamente mediante el análisis de sensibilidad de la sección siguiente.

\section{Análisis de Sensibilidad de la Función de Recompensa}
\label{sec:sensitivity}

Para evaluar la robustez de la calibración y justificar la elección del perfil principal, se ha entrenado un agente PPO bajo cada uno de los cuatro perfiles, manteniendo idénticos el resto de hiperparámetros (200~000 pasos por configuración, partición \textit{train}/\textit{test} 80/20). Los resultados sobre el periodo \textit{out-of-sample} se resumen en la tabla \ref{tab:sensitivity}.

\begin{table}[h]
\centering
\renewcommand{\arraystretch}{1.25}
\begin{tabular}{|l|c|c|c|c|c|}
\hline
\textbf{Configuración} & $\boldsymbol{\varphi}$ & $\boldsymbol{\gamma}$ & \textbf{Sharpe} & \textbf{Retorno (\%)} & \textbf{MDD (\%)} \\ \hline
A: \texttt{balanced} & 0.02 & 0.010 & 1.368 & 112.4 & $-31.8$ \\ \hline
B: \texttt{conservative} & 0.05 & 0.010 & 1.613 & 170.6 & $-40.8$ \\ \hline
\textbf{C: \texttt{low\_turnover}} \textbf{(TFM)} & \textbf{0.02} & \textbf{0.020} & \textbf{1.770} & \textbf{182.2} & $-40.0$ \\ \hline
D: \texttt{aggressive} & 0.01 & 0.005 & 1.310 & 102.6 & $-31.1$ \\ \hline
\end{tabular}
\caption{Resultados del análisis de sensibilidad sobre el periodo out-of-sample.}
\label{tab:sensitivity}
\end{table}

La lectura de la tabla \ref{tab:sensitivity} permite extraer tres conclusiones:

\begin{enumerate}
    \item \textbf{\texttt{low\_turnover} maximiza Sharpe y retorno.} Con un Sharpe de 1.77 y un retorno acumulado del 182\%, supera a las otras tres configuraciones. Este resultado justifica empíricamente —no narrativamente— su elección como perfil principal.
    \item \textbf{Trade-off MDD vs.\ Sharpe.} Las configuraciones \texttt{balanced} y \texttt{aggressive} obtienen menor MDD ($\sim$32\%) frente al $\sim$40\% de \texttt{conservative} y \texttt{low\_turnover}. La elección de \texttt{low\_turnover} prioriza el equilibrio retorno/riesgo capturado por Sharpe, asumiendo que (i) el propio $\gamma$ ya penaliza la rotación durante el entrenamiento y (ii) el Sharpe ya pondera retorno por volatilidad.
    \item \textbf{Robustez frente a la calibración.} El rango de Sharpe (1.31 a 1.77) es estrecho frente a la magnitud de las variaciones de $\varphi$ y $\gamma$. Esto sugiere que la política aprendida no es excesivamente sensible a la calibración exacta de la recompensa dentro del rango explorado.
\end{enumerate}

\section{Agente PPO: Arquitectura, Hiperparámetros e Inferencia}
\label{sec:agente_ppo}

\subsection{Arquitectura Actor-Crítico}

El agente se implementa con la API \texttt{PPO} de \textit{Stable-Baselines3}, instanciada con la política \texttt{MlpPolicy}, que corresponde a la \texttt{ActorCriticPolicy} estándar de la librería. Esta política define dos cabezas sobre una red MLP compartida:

\begin{itemize}
    \item Un \textit{actor} que produce los parámetros de una distribución Gaussiana diagonal sobre el espacio de acciones $[0,1]^{9}$.
    \item Un \textit{crítico} que estima la función de valor $V_{\phi}(\mathbf{s})$, utilizada para calcular ventajas $\hat{A}_t$ con baja varianza mediante el algoritmo \textit{Generalized Advantage Estimation} (GAE) estándar de SB3.
\end{itemize}

La arquitectura concreta es un MLP de dos capas ocultas de 256 unidades cada una (\texttt{net\_arch=[256, 256]}), valor por defecto recomendado por SB3 para problemas continuos de dimensionalidad media.

\subsection{Hiperparámetros}

La tabla \ref{tab:ppo_params} resume los hiperparámetros utilizados, idénticos en producción y en validación temporal. La unificación entre fases es deliberada: garantiza que las métricas reportadas reflejan el desempeño esperable del modelo desplegado y no de una versión simplificada para reducir cómputo.

\begin{table}[h]
\centering
\renewcommand{\arraystretch}{1.25}
\begin{tabular}{|l|l|p{6.0cm}|}
\hline
\textbf{Hiperparámetro} & \textbf{Valor} & \textbf{Justificación} \\ \hline
Política & \texttt{MlpPolicy} & Actor-crítico estándar de SB3 con MLP compartido. \\ \hline
\texttt{net\_arch} & $[256, 256]$ & Capacidad razonable para la dimensionalidad del estado. \\ \hline
\texttt{learning\_rate} & $1\!\times\!10^{-4}$ & Actualizaciones más graduales que el valor por defecto $3\!\times\!10^{-4}$, útil en señales con baja relación señal-ruido. \\ \hline
\texttt{n\_steps} & 2048 & Trayectoria larga entre actualizaciones; reduce la varianza de las estimaciones de ventaja. \\ \hline
\texttt{batch\_size} & 128 & Gradientes más estables que con \texttt{batch\_size} 64. \\ \hline
\texttt{clip\_range} ($\epsilon$) & 0.10 & Recorte conservador del PPO; limita los pasos en una señal con baja relación señal-ruido. \\ \hline
\texttt{ent\_coef} & 0.01 & Bonificación por entropía para mantener exploración a lo largo del entrenamiento. \\ \hline
\texttt{vf\_coef} & 0.5 & Peso del error del crítico en la pérdida total (valor por defecto). \\ \hline
\texttt{max\_grad\_norm} & 0.5 & \textit{Clipping} de gradientes para evitar inestabilidades. \\ \hline
\end{tabular}
\caption{Hiperparámetros del entrenamiento PPO.}
\label{tab:ppo_params}
\end{table}

\subsection{Modo de Inferencia: Determinismo en el Backtest}
\label{sec:inferencia_determinista}

PPO aprende una política estocástica $\pi_{\theta}(\mathbf{a} \mid \mathbf{s})$, modelada en este caso como una Gaussiana diagonal sobre el simplex. Durante el entrenamiento, esta estocasticidad es imprescindible para explorar el espacio de asignaciones. En cambio, en todas las fases de inferencia —\textit{walk-forward}, \textit{expanding window}, análisis de sensibilidad, análisis por régimen y simulaciones del usuario final— se utiliza el modo determinista (\texttt{deterministic=True} en la API SB3), que devuelve la media de la distribución del actor en lugar de muestrear. Esta decisión se justifica por tres razones:

\begin{itemize}
    \item \textbf{Reproducibilidad.} Dos ejecuciones del mismo modelo sobre los mismos datos producen exactamente la misma curva de \textit{equity}, lo que permite comparar con rigor entre configuraciones, perfiles y \textit{baselines}.
    \item \textbf{Realismo operativo.} Un gestor patrimonial requiere la mejor estimación que la política aprendida puede ofrecer, no una muestra aleatoria. La media de la distribución corresponde a la asignación modal de la política.
    \item \textbf{Separación entre entrenamiento y evaluación.} La estocasticidad solo es necesaria durante el entrenamiento; mantenerla en evaluación contaminaría las métricas con ruido ajeno al desempeño real.
\end{itemize}

\section{Baselines del Estudio Comparativo}
\label{sec:baselines}

Para evaluar el agente PPO se han implementado seis estrategias de referencia (\textit{baselines}) que cubren un espectro deliberadamente amplio: desde la cordura mínima exigible (\textit{lower bound}) hasta \textit{baselines} competitivas del sector cuantitativo. La tabla \ref{tab:baselines} resume el catálogo.

\begin{table}[h]
\centering
\renewcommand{\arraystretch}{1.25}
\begin{tabular}{|l|p{8.5cm}|p{2.0cm}|}
\hline
\textbf{Baseline} & \textbf{Descripción y rol} & \textbf{Rebalanceo} \\ \hline
Equal-Weight & 1/N a cada activo del universo. \textit{Baseline} pasiva clásica. & Mensual \\ \hline
60/40 & 60\% IVV (renta variable) + 40\% BND (renta fija). Cartera institucional canónica. & Mensual \\ \hline
Buy \& Hold & 1/N inicial sin rebalanceo posterior. \textit{Baseline} de mínima fricción. & Ninguno \\ \hline
Markowitz MV & Optimización media-varianza con ventana de estimación de 252 días, maximizando Sharpe vía SLSQP. \textit{Baseline} teórica clásica \cite{markowitz1952}. & Mensual \\ \hline
Random Uniform & Pesos muestreados uniformemente del simplex (Dirichlet$(\mathbf{1})$) con semilla fija. \textit{Lower bound} de cordura: si el agente DRL no supera a esta \textit{baseline}, no está aprendiendo. & Mensual \\ \hline
Momentum Top-K & Cada mes, asigna pesos equiponderados al \textit{top}-3 de activos con mejor retorno acumulado en los últimos 60 días. Implementa el factor \textit{momentum} \textit{cross-sectional} documentado por \cite{jegadeesh1993}. \textit{Baseline} competitiva del sector cuantitativo. & Mensual \\ \hline
\end{tabular}
\caption{Estrategias de referencia (baselines) implementadas.}
\label{tab:baselines}
\end{table}

La inclusión de Random Uniform y Momentum Top-K responde a dos preguntas distintas: la primera verifica que el agente realmente aprende y no replica el azar; la segunda verifica que el agente aporta valor por encima de un factor financiero conocido y operativamente sencillo. Si el PPO no supera al Momentum Top-K, su justificación como artefacto del TFM se debilita; si sí lo hace consistentemente, se valida la hipótesis de que un agente DRL aporta información adicional sobre la simple persistencia del momentum.

\section{Protocolo de Validación Temporal}
\label{sec:validacion}

La validación del agente sigue dos esquemas complementarios habituales en la literatura de \textit{Machine Learning} financiero \cite{lopezdeprado2018}:

\begin{itemize}
    \item \textbf{\textit{Walk-forward} (rolling).} Se desliza una ventana temporal fija a lo largo del histórico: el modelo se entrena en un bloque \textit{train} y se evalúa inmediatamente después en un bloque \textit{test} contiguo, repitiendo el procedimiento sucesivamente. Cada ventana proporciona una estimación de desempeño \textit{out-of-sample} en condiciones realistas, evaluando exclusivamente sobre datos posteriores al entrenamiento.
    \item \textbf{\textit{Expanding window}.} Variante en la que el bloque \textit{train} crece progresivamente con el tiempo (acumula toda la historia disponible hasta el corte), mientras que el bloque \textit{test} se desliza sobre el periodo siguiente. Este esquema simula mejor la situación de un sistema en producción que reentrena periódicamente con todos los datos acumulados.
\end{itemize}

En ambos esquemas, los hiperparámetros del entrenamiento son idénticos a los de la configuración de producción (tabla \ref{tab:ppo_params}), de modo que las métricas reportadas reflejan el desempeño esperable del modelo desplegado y no de una versión simplificada. Los resultados de cada ventana se agregan en sus medias y desviaciones, y se reportan junto con las correspondientes curvas de \textit{equity} comparadas con las seis \textit{baselines} de la sección \ref{sec:baselines}.

\chapter{Glosario}

\textit{Deep Reinforcement Learning} (DRL), \textit{Asset Allocation}, ETFs de Criptomonedas, Optimización de Carteras, PPO (\textit{Proximal Policy Optimization}), Finanzas Cuantitativas.



\clearpage
\phantomsection
\renewcommand{\bibname}{Bibliografia}
\addcontentsline{toc}{chapter}{Bibliografia}

\bibliographystyle{unsrtnat}
\bibliography{bibliografia}


\begin{comment}
\chapter{Bibliografía}

\begin{itemize}
    \item \textbf{Sutton, R. S., \& Barto, A. G. (2018).} \textit{Reinforcement Learning: An Introduction}. MIT Press. Fundamentos teóricos del RL.
    \item \textbf{Lopez de Prado, M. (2018).} \textit{Advances in Financial Machine Learning}. Wiley. Metodología de validación \textit{out-of-sample}.
    \item \textbf{Arulkumaran, K., et al. (2017).} \textit{Deep Reinforcement Learning: A Brief Survey}. IEEE Signal Processing Magazine. Justificación de los avances en Deep RL.
    \item \textbf{Liu, X. Y., et al. (2021).} \textit{FinRL: A Deep Reinforcement Learning Library for Automated Stock Trading}. .
    \item \textbf{Schulman, J., et al. (2017).} \textit{Proximal Policy Optimization Algorithms}. OpenAI. 
    \item \textbf{Markowitz, H. (1952).} \textit{Portfolio Selection}.

    \item \textbf{Liu, X. Y., et al. (2018).} \textit{FinRL: A Deep Reinforcement Learning Library}.
    \item \textbf{Hougan, M., \& Lawant, D. (2021).} \textit{Cryptoassets}.
    \item \textbf{Oates, B. J. (2006).} \textit{Researching Information Systems and Computing}.
\end{itemize}
\end{comment}
\newpage
\appendix


\end{document}
